% titel: Blatt 7, A6: Umkehrregel
% skript: Kap. 4.1 (Umkehrfunktion, Umkehrregel)

% Teil I: Werkzeuge, kein Lösungsweg.

Für die Ableitung einer Umkehrfunktion braucht man die Umkehrregel. Sie
verbindet $f'$ und $(f^{-1})'$ — dort, wo $f'$ nicht verschwindet und die
Umkehrfunktion differenzierbar ist.

\begin{werkzeugkasten}[Umkehrregel]
Ist $f$ auf einem Intervall streng monoton und differenzierbar mit $f'\neq0$
auf dem Intervall, so ist $f^{-1}$ differenzierbar und
\begin{equation}
  (f^{-1})'(y)=\frac{1}{f'(f^{-1}(y))}
  \qquad\text{bzw.}\qquad
  \dv{}{y}f^{-1}(y)=\frac{1}{f'(x)},\quad x=f^{-1}(y).
\end{equation}
Man leitet die Identität $f^{-1}(f(x))=x$ nach $x$ ab und löst auf; das ist
der Kern des Beweises.

\textit{Skript:} Kap. 4.1, Umkehrregel.
\end{werkzeugkasten}

\begin{werkzeugkasten}[Kosinus und arccos]
$\cos$ ist auf $[0,\pi]$ streng monoton fallend und bildet $[0,\pi]$ auf
$[-1,1]$ ab; die Umkehrfunktion
$\arccos:[-1,1]\to[0,\pi]$ ist die Umkehrung. Für $|y|<1$ gilt die
Pythagoras-Beziehung im Einheitskreis: ist $x=\arccos y$, so ist $\cos x=y$
und $\sin x=\sqrt{1-y^2}$.

\textit{Skript:} Kap. 2.4 (Umkehrfunktionen) und Kap. 4.1.
\end{werkzeugkasten}

\begin{achtung}[Darauf achten]
\begin{itemize}
  \item $\cos'=-\sin$ — das Minuszeichen nicht vergessen; es bewirkt das
        Minus in der arccos-Ableitung.
  \item Die Umkehrregel versagt, wo $f'=0$ ist. Hier: $\sin x=0$ bei
        $x=0,\pi$, also $y=\pm1$; dort ist arccos nicht differenzierbar
        (der Nenner $\sqrt{1-y^2}$ wird null).
  \item Argument und Definitionsbereich: arccos ist nur auf $[-1,1]$ definiert;
        die Formel gilt nur für $|y|<1$.
\end{itemize}
\end{achtung}
